\clearpage
\begin{center}\thispagestyle{empty}
{\LARGE\bfseries{概要}}
\end{center}

大規模言語モデル（LLM）の複数体を議論させるマルチエージェント推論は，単一モデルの限界を超える手段として注目されるが，小型モデルでは「計算量を揃えた多数決（Self-Consistency）に勝てない」という強い批判がある．本研究は，4B級モデル（Qwen3-4B-Instruct-2507）にLoRAで異なるペルソナを焼き込んだ3体のエージェント集団を構成し，（1）議論によるチーム推論が生成回数を揃えたSelf-Consistency@9（SC@9）を上回れるか，（2）チームレベル適応度（厳密Shapley値$\times$fitness sharing）による重み空間の進化的最適化が寄与するかを，3ベンチマーク（MMLU-Pro・MATH-500・SuperGPQA）$\times$3シードの対応あり検定（問題IDクラスタリングを考慮した置換検定・ブートストラップ＋Holm補正）で検証した．

実験1（主実験）の結果，議論の効果はベンチマーク依存に反転することが判明した．素の議論は数学（$+4.2$pt）と高難度科学（$+3.1$pt）で有意に有効だがMMLU-Proでは有意に有害（$-4.1$pt）であり，重みペルソナチームは逆にMMLU-Proのみで有効（対ベース$+4.5$pt）だった．この反転は「多様性の利得」と「ペルソナ注入による能力毀損」の2軸トレードオフとして整理できる（ただし反転を問題の検証可能性に帰する解釈はベンチマーク内のカテゴリ別分析では支持されず，仮説にとどまる）．診断分析により，（i）ペルソナSFTが思考連鎖を半減させ長い導出能力を選択的に破壊すること，（ii）3体の誰かが正解している問題の最大13\%分を多数決集約が取りこぼすこと，（iii）進化は適応度測定ノイズ（100問で$\pm 4.9$pt）下の勝者の呪いにより選抜として機能しなかったうえ，進化前後の$\Delta W$がcosine類似度0.9997以上とほぼ同一——すなわち進化演算子が重み空間をほとんど動かしていなかったことを明らかにした．

実験2では診断に基づく処方——ベースモデル自己生成の検証済み長CoTリプレイ混合（38\%）による再学習，自己誤り特定時のみ回答変更を許す条件付き議論，発言匿名化，logprob重み付き投票——を，事前定義した合格基準を持つ実測ゲート方式で検証した．この過程で評価コンテナイメージ間の系統差（同一問題・同一設定で$+6$pt級，ベースモデルの素の生成に選択的に作用）を検出したため，主要4条件を同一環境で全面再測定して統制した．最終評価において，処方はペルソナSFTが生んだ能力毀損をベース単体との同等性（等価マージン$\pm 2$pt）が成立する水準まで回復させ（対旧チーム$+2.7$pt，$p=0.0003$；MATH-500では対ベース$+3.4$pt有意優位），一方でSC@9には$-3.5$pt（$p<10^{-4}$）と有意に及ばなかった．

本研究の貢献は，（a）小型LLM議論チームの成否を2軸トレードオフと損失分解で説明する診断枠組み，（b）進化ループの機能不全（探索・多様性圧・選抜の三重の不全）の実測による特定と選抜ノイズの定量化，（c）能力毀損を回復させる処方の実証と，それを尽くしてもなお生成回数を揃えた多数サンプリングに及ばないという統制された負の結果，（d）評価環境の系統差がベースライン側を選択的に過小評価しチーム手法を見かけ上優位にするバイアスの発見，である．数学（MATH-500）では議論チームがベース単体を有意に上回り続けており，マルチエージェント議論の価値が残る条件を実測に基づいて画定した．




\clearpage
\phantom{s}\thispagestyle{empty}
\clearpage

\pagestyle{empty}
